\documentclass[11pt]{article}
\usepackage[margin=1in]{geometry}
\usepackage{booktabs}
\usepackage{graphicx}
\usepackage[hidelinks]{hyperref}
\usepackage{enumitem}
\usepackage{titlesec}
\titleformat{\section}{\normalfont\large\bfseries}{\thesection.}{0.6em}{}
\titleformat{\subsection}{\normalfont\normalsize\bfseries}{\thesubsection}{0.6em}{}

\begin{document}

\begin{center}
{\LARGE\bfseries Auditing Epistemic Guardrails From Their Own Artifacts:\\[2pt]
Independent Verification of Subjective-Logic Uncertainty Claims\\ for LLM-Driven Actuation\par}
\vspace{10pt}
Muntaser Syed\\
Florida Institute of Technology\\
\texttt{msyed2011@my.fit.edu}\\[4pt]
{\small With Apart Research\footnote{Research conducted at the AI x Epistemics Research Sprint, November 2026.}}
\end{center}

\vspace{4pt}
\noindent\textbf{Abstract.}
Claims about epistemic guardrails are becoming load-bearing: if an uncertainty layer is said to be what stands between a language model and a physical actuator, other researchers and practitioners will build on that claim. Yet such claims are rarely re-derived from the artifacts released with them. This report is an artifact audit of one published system, \textit{Epistemic Edge} (Syed and Silaghi, IEEE IRI 2026), which places subjective-logic uncertainty quantification between IoT sensors and an LLM that can actuate. I recomputed the paper's three headline claims directly from its released ground-truth result files, without re-running any model: all 17 released ablation runs (6,800 trials, including runs outside the published seven-model subset), and the BATADAL benchmark summaries, configuration heatmap, and per-timestamp scores (422 configurations). The central claim reproduces exactly: threshold-guardrail accuracy is 1.00 with the full pipeline and exactly 0.40 without the epistemic layer, in 17 of 17 runs. The BATADAL AUROC of 0.9004 re-derives exactly, and the audit localizes the failure mode precisely: accuracy collapses only when the fused opinion is vacuous or absent, not when fusion is bypassed by passthrough. Per-event analysis confirms that no single subjective-logic signal detects all five attacks best. Artifact-complete releases make claims checkable; this audit shows what checking them adds.

\section{Introduction}

AI systems are increasingly asked not only to describe the world but to act on it, and to say how sure they are. The sprint's first two tracks ask the same question from opposite ends: does a model know how solid its claims are, and what trust infrastructure makes checking claims cheap and hard to manipulate? A precondition for both is that published epistemic claims are themselves checkable. In practice they usually are not: papers report aggregate numbers, and the result files behind them, if released at all, are seldom re-examined.

This project takes one published epistemic-guardrail system and audits its claims against its own released artifacts. The system is \textit{Epistemic Edge}~\cite{syed2026epistemic}, a four-tier neuro-symbolic pipeline that annotates sensor streams with subjective-logic (SL) opinions (belief, disbelief, uncertainty)~\cite{josang2016}, fuses them, decays them over time, and lets an LLM actuate only when the fused opinion passes an epistemic threshold guardrail in addition to a behavioral whitelist. The paper reports three findings: (1) removing the epistemic layer collapses threshold-guardrail accuracy from 1.00 to exactly 0.40 across seven locally deployed LLMs (2,800 controlled trials); (2) on the BATADAL water-distribution attack benchmark~\cite{taormina2018}, calibrated SL opinions reach AUROC 0.9004 with no trained weights; (3) the SL dimensions are decomposable: different attack events are optimally detected by different SL-derived signals, and no single signal wins on all five events.

\textbf{Disclosure and scope.} The audited paper is my own prior published work with Marius Silaghi, and the audited artifacts are its public release. This is therefore a self-audit, and Section 5 states what that does and does not establish. The sprint-period work reported here is the audit itself: the recomputation of every headline number from the released files, two analyses the paper does not contain (a paired fusion analysis over the full configuration grid, and a per-event signal analysis under a stated protocol), and this report. The system, the paper, and the underlying experiment data are preexisting and are marked as such throughout.

Our main contributions are:
\begin{enumerate}[leftmargin=*,itemsep=1pt]
\item A complete artifact audit of a published epistemic-guardrail claim set: all 17 released ablation runs (6,800 trials) and all 422 BATADAL configurations, recomputed with a single released script, rather than the published subsets alone.
\item An exact verification of the two quantitative headline claims (1.00 vs.\ 0.40; AUROC 0.9004), plus a sharper localization of the failure mode: the collapse occurs exactly when the fused opinion is vacuous or absent, and not under passthrough fusion.
\item A per-event re-analysis showing the decomposability claim holds in substance (three different signals win the five events; belief-family signals are anti-predictive in the positive direction), and an audit observation about phase-dependent configuration selection that headline reading alone would miss.
\end{enumerate}

\section{Related Work}

Subjective logic~\cite{josang2016} represents an opinion as a triple of belief, disbelief, and uncertainty summing to one, with operators for fusion and discounting; it is one of several formalisms for carrying explicit uncertainty alongside a decision. The broader calibration literature shows that modern neural networks are typically miscalibrated out of the box~\cite{guo2017}, and recent work argues that common benchmarks reward confident guessing over honest abstention~\cite{kalai2025}, while benchmarks such as MASK measure whether models contradict their own stated beliefs under pressure~\cite{ren2025}. Epistemic Edge sits downstream of these questions: it does not ask the LLM to be calibrated, but wraps it in an external epistemic layer whose opinions are computed from sensor evidence, and blocks actuation when uncertainty is high.

BATADAL~\cite{taormina2018} is a standard benchmark for cyber-attack detection in water distribution systems, with labeled attack events in real telemetry. Artifact and reproducibility evaluation has a long history in systems venues, but remains rare for safety-relevant ML claims, where the cost of an unchecked number is highest. This audit is a small instance of a general practice the field needs: claims shipped with complete, immutable result files, and independent recomputation treated as normal scientific work.

\section{Methods}

\textbf{Artifacts audited.} The public repository\footnote{\url{https://github.com/jemsbhai/epistemic-edge}} (also on PyPI as \texttt{epistemic-edge}) declares its \texttt{experiments/results/} JSON and CSV files to be the immutable ground truth for every reported number. The audit covers: 17 ablation result files for the controlled-trial experiment; 3 BATADAL summary files; the primary run's 422-row configuration heatmap; the primary run's per-timestamp score series (4,177 evaluation rows, 43 sensors); and the run's reproducibility record.

\textbf{Audit protocol.} For each headline claim I extracted the corresponding quantity from every released file that records it, not only from the files behind the published tables, and compared values exactly. No LLM was re-run and the BATADAL pipeline was not re-executed: the audit verifies that the released record supports the published numbers, and says so where verification of a different kind (for example, re-running the trials) would be needed. The audit script is released with this report\footnote{\url{https://github.com/jemsbhai/epistemic-edge-audit}} and uses only the Python standard library.

\textbf{Conditions.} The controlled-trial experiment runs eight conditions per model over five IoT actuation scenarios, 400 trials per run. Threshold-guardrail accuracy is defined for the five conditions in which the threshold guardrail is active: C (full pipeline), D (no temporal decay), E1 (no fusion, vacuous opinion), E2 (no fusion, passthrough opinion), and G (no epistemic layer).

\textbf{Audit-side analyses.} Two analyses are new. (i) A paired fusion analysis: across the 422-configuration BATADAL grid, I matched configurations differing only in whether SL fusion is enabled (150 matched pairs) and compared headline AUROC. (ii) A per-event analysis: using the released per-timestamp scores of the primary run, I computed, for each of the five attack events, the rank-based AUROC of each SL-derived signal with event-window samples as positives and all non-attack samples as negatives. The paper's own per-event analysis used its authors' protocol; mine is stated here so the comparison is honest.

\section{Results}

\subsection{Claim 1: the epistemic layer is necessary (verified exactly, and beyond the published subset)}

Table~\ref{tab:ablation} summarizes the audit of all 17 released ablation runs. In every run, the full pipeline scores threshold-guardrail accuracy 1.00 and the no-epistemic condition scores exactly 0.40, matching the paper's claim for its published seven-model subset (2,800 trials) and extending it to the full released record (6,800 trials, including repeat and variant runs the paper does not tabulate).

The condition structure also sharpens the claim. The 0.40 floor appears in exactly two conditions: G (no epistemic layer) and E1 (fusion replaced by the vacuous opinion). Condition E2, which keeps per-source opinions but passes them through without fusion, retains 1.00, as does D (no decay). The collapse is therefore not caused by the absence of fusion machinery as such; it occurs when the opinion the guardrail sees is vacuous or missing, which is precisely the state in which a threshold on uncertainty cannot discriminate. The active ingredient, on this benchmark, is a non-vacuous calibrated opinion at the guardrail.

\begin{table}[t]
\centering
\small
\begin{tabular}{@{}llll@{}}
\toprule
Condition & Description & Threshold accuracy & Runs \\
\midrule
C & Full pipeline & 1.00 & 17/17 \\
D & No temporal decay & 1.00 & 17/17 \\
E2 & No fusion, passthrough & 1.00 & 17/17 \\
E1 & No fusion, vacuous opinion & 0.40 & 17/17 \\
G & No epistemic layer & 0.40 & 17/17 \\
\midrule
\multicolumn{4}{@{}l}{\footnotesize 17 released runs, 400 trials each (6,800 total). A (bare LLM), B (raw payloads),}\\
\multicolumn{4}{@{}l}{\footnotesize and F (no guardrails) do not define threshold accuracy.}\\
\bottomrule
\end{tabular}
\caption{Artifact audit of Claim 1. Values recomputed from every released ablation result file.}
\label{tab:ablation}
\end{table}

\subsection{Claim 2: BATADAL AUROC 0.9004 (verified exactly)}

The final BATADAL summary records a best AUROC of 0.9004, attained by the configuration: historical calibration strategy, window $k=24$, decay $\lambda=0.25$, fusion and decay enabled, domain physics bounds, fused-disbelief anomaly score, smoothing 5. The run's reproducibility record (seed 42) reports original and verification metrics identical on every field (AUROC 0.9004, AUPRC 0.6379, best F1 0.6481), and my recomputation from the summary agrees. Two earlier summary files record 0.8447, which the audit traces to an earlier selection stage: 0.8447 is exactly the grid maximum of the fused-uncertainty signal (Section 4.3). The headline number is the final configuration's, not the largest number anywhere in the record.

The paired fusion analysis supports the paper's reported fusion effect in direction: across 150 matched configuration pairs differing only in fusion, AUROC improves by a mean of $+0.065$ (median $+0.054$); 65.3\% of pairs improve and 33.3\% are exactly unchanged (fusion is a no-op where a timestamp carries effectively one opinion stream). I did not re-derive the paper's Wilcoxon statistic ($p<10^{-4}$, Cohen's $d=1.01$) on its original subset; the paired analysis here is an independent check on the full grid, and it agrees in direction and rough magnitude.

One audit observation: the full grid contains a fused-disbelief AUROC of 0.9356, higher than the selected configuration's 0.9004. It occurs in a decay-ablation-phase row (decay off, $\lambda=0.1$, smoothing 1) whose selected anomaly score was a different signal; within the signal-sweep phase in which selection was performed, the selected row is the best row. The published number is thus phase-consistent, but the record shows why artifact access matters: a reader of only the heatmap could reasonably ask why 0.9356 was not the headline, and only the phase structure answers it.

\subsection{Claim 3: signals are decomposable (verified in substance)}

Across the 422-configuration grid, the per-signal AUROC maxima differ widely: fused disbelief 0.9356, composite uncertainty-disbelief 0.9322, max disbelief 0.8764, fused uncertainty 0.8447, mean uncertainty 0.8151, max uncertainty 0.7347. No signal's maximum is attained by the same configuration as another's, and the ranking of signals changes across the grid.

The per-event re-analysis (Table~\ref{tab:perevent}) is more direct. Three different signals win the five events: fused uncertainty wins events 1, 4, and 5 (0.995, 0.887, 0.996), fused disbelief wins event 2 (0.805), and max disbelief wins event 3 (0.666). Event 3 is the hardest for every signal. Belief-family signals are strongly anti-predictive as positive-direction scores (fused belief: 0.008 and 0.004 on events 1 and 5), which is the expected signature of belief collapsing during an attack: read in the negative direction they are among the most informative signals on those events. The paper's claim, that no single signal wins on all five events and that the SL dimensions detect structurally different attack classes, holds under this independent protocol; the stronger phrasing that every event has a distinct optimal signal does not (fused uncertainty wins three of five).

\begin{table}[t]
\centering
\small
\begin{tabular}{@{}lcccc@{}}
\toprule
Event & Rows & Best signal & AUROC & Runner-up \\
\midrule
1 & 1731--1772 & fused uncertainty & 0.995 & fused disbelief 0.992 \\
2 & 2337--2396 & fused disbelief & 0.805 & fused uncertainty 0.750 \\
3 & 2851--2887 & max disbelief & 0.666 & fused disbelief 0.616 \\
4 & 3508--3514 & fused uncertainty & 0.887 & fused disbelief 0.854 \\
5 & 3741--3813 & fused uncertainty & 0.996 & fused disbelief 0.996 \\
\bottomrule
\end{tabular}
\caption{Per-event re-analysis from the released primary-run score series, under the audit protocol of Section 3. Event windows match the released summary's attack records exactly.}
\label{tab:perevent}
\end{table}

\section{Discussion and Limitations}

The audit's positive result is narrow and worth stating narrowly: the published numbers of Epistemic Edge are exactly the numbers its released artifacts contain, across a record roughly 2.4 times larger than the published tables, and the one apparent anomaly in the record (a grid AUROC above the headline) has a documented structural explanation. That is the outcome trust infrastructure is supposed to produce, and it was cheap: one script, no compute of consequence, because the release was artifact-complete.

\textbf{Limitations.} First, this is a self-audit: I am the first author of the audited paper. The mitigation is procedural, not personal: every step is recomputation from public immutable files with a released script, so any reader can repeat it exactly; but an auditor with no stake in the outcome would be stronger evidence, and this report should be read as an invitation to replicate rather than as independent certification. Second, the audit verifies the record, not the original runs: a result file fabricated at the source would pass. Provenance of the files themselves (for example, signed or timestamped result logs) is outside what the release provides, and is a concrete improvement future releases should make. Third, the controlled-trial benchmark is small and synthetic (five scenarios); the 1.00/0.40 separation is exact on it, and says little about effect sizes on messier actuation tasks. Fourth, my per-event protocol differs from the paper's, so Claim 3 is verified in substance rather than to the digit.

\textbf{Future work.} Three extensions follow naturally: third-party replication of the ablation on independently run trials; provenance-signed result artifacts in the epistemic-edge release, so future audits can check origin as well as content; and applying the same audit treatment to the system's calibration claims under distribution shift, where the guardrail matters most and the current record is thinnest.

\section{Conclusion}

Epistemic claims deserve the same skepticism they are built to operationalize. Auditing one published guardrail system against its own artifacts, I found its headline claims exact, its failure mode more precisely located than the paper states (vacuous or absent opinions, not missing fusion), and its decomposability claim supported in substance. The larger point is procedural: none of this required new experiments, only a complete release and the habit of checking. Benchmarks and guardrails will be trusted in proportion to how checkable they are.

\section*{Code and Data}
\begin{itemize}[leftmargin=*,itemsep=0pt]
\item Audited system and artifacts: \url{https://github.com/jemsbhai/epistemic-edge} (PyPI: \texttt{epistemic-edge}). Preexisting work, published at IEEE IRI 2026.
\item This audit (script, full output, report source): \url{https://github.com/jemsbhai/epistemic-edge-audit}. Sprint-period work.
\item Dataset: BATADAL (datasets 03/04), distributed by the BATADAL organizers; the released summaries, heatmap, and score series used here are in the audited repository.
\end{itemize}

\section*{Author Contributions}
Muntaser Syed designed and executed the audit, wrote the audit script and analyses, and wrote this report. The audited system and paper are prior joint work with Marius Silaghi (Florida Institute of Technology), who is not an author of this sprint report.

\begin{thebibliography}{9}\small
\bibitem{syed2026epistemic} M. Syed and M. C. Silaghi. Epistemic Edge: Subjective Logic Guardrails for LLM-Driven IoT Actuation. \textit{Proc.\ IEEE 27th International Conference on Information Reuse and Integration for Data Science (IRI)}, 2026, pp. 37--40.
\bibitem{josang2016} A. J\o sang. \textit{Subjective Logic: A Formalism for Reasoning Under Uncertainty}. Springer, 2016.
\bibitem{taormina2018} R. Taormina et al. Battle of the Attack Detection Algorithms: Disclosing Cyber Threats on Water Distribution Systems. \textit{Journal of Water Resources Planning and Management}, 144(8), 2018.
\bibitem{guo2017} C. Guo, G. Pleiss, Y. Sun, and K. Q. Weinberger. On Calibration of Modern Neural Networks. \textit{Proc.\ ICML}, 2017.
\bibitem{kalai2025} A. T. Kalai, O. Nachum, S. S. Vempala, and E. Zhang. Why Language Models Hallucinate. arXiv preprint, 2025.
\bibitem{ren2025} J. Ren et al. The MASK Benchmark: Disentangling Honesty From Accuracy in AI Systems. arXiv preprint, 2025.
\end{thebibliography}

\appendix
\section{Limitations and Dual-Use / Ethical Considerations}
\textbf{What this result does and does not establish.} The audit establishes that one paper's published numbers match its released record. It does not establish that the underlying system is safe to deploy, that its guardrails generalize beyond the tested scenarios, or that other systems' claims would survive the same treatment. The self-audit limitation of Section 5 applies in full.

\textbf{Dual use.} Localizing a guardrail's failure mode (vacuous or absent opinions at the guardrail) could in principle help an attacker craft sensor inputs that drive a deployed system's fused opinion toward vacuity. Three considerations bound this risk: the audited system is published and fully open, so the ablation structure was already public; the audit introduces no new attack method and tests no live system; and the defensive value of verification practice, knowing precisely which component carries the safety claim, is the precondition for hardening it. On balance the audit is a defensive contribution.

\textbf{Data and human subjects.} No personal data was collected or used. BATADAL is a public benchmark of water-utility telemetry released for attack-detection research; no human subjects are involved.

\section{LLM Usage Statement}
An AI assistant was used, under the author's direction, to write and run the audit scripts and to assist in drafting this report. All numbers reported here were recomputed from the released artifacts by the audit script in the project repository, and every claim was checked against the script's output before inclusion. The author is responsible for all content.

\end{document}
